\uuid{AX5X}
\exo7id{3878}
\titre{exo7 3878}
\auteur{quercia}
\organisation{exo7}
\datecreate{2010-03-11}
\isIndication{false}
\isCorrection{false}
\chapitre{Continuité, limite et étude de fonctions réelles}
\sousChapitre{Etude de fonctions}
\cpge{mpsi-mp2i.limites-continuite.a, pcsi.limites-continuite.a, ptsi.limites-continuite.a, mpsi-mp2i.limites-continuite.c, pcsi.limites-continuite.c, ptsi.limites-continuite.c, mpsi-mp2i.denombrement.a, pcsi.denombrement.a, ptsi.denombrement.a, bcpst1.limites-continuite.a, bcpst1.limites-continuite.c, bcpst1.denombrement, mp-mpi.variables-aleatoires-discretes.a, pc.familles-sommables, psi.denombrables-familles-sommables, pt.probabilites-lois-usuelles.a}
\module{Analyse}
\niveau{L1}
\difficulte{}

\contenu{
\texte{
Soit $f : {[a,b]} \to {\R}$ monotone.
Pour $x \in {]a,b[}$,
on pose $\delta(x) = \bigl|\lim_{y\to x^+} f(y) - \lim_{y\to x^-} f(y)\bigr|$
(saut de $f$ en $x$).
}
\begin{enumerate}
    \item \question{Pour $n \in \N^*$, montrer que
      $E_n = \left\{ x \in {]a,b[} \text{ tq } \delta(x) > \frac 1n \right\}$ est fini.}
    \item \question{En déduire que l'ensemble des points de discontinuité de $f$ est au plus dénombrable.}
\end{enumerate}
}
